\begin{sgastatement}{\SGAInlineRefTarget{sga3:xxvi:corollary:4-2-5}Corollary 4.2.5}
Let \(S\) be a scheme, \(G\) an \(S\)-reductive group, and \(P,Q\) two
parabolic subgroups of \(G\), with \((P,Q)\) in transversal position.
\begin{enumerate}[label=\textup{(\roman*)}]
\item Let \(P',Q'\) be parabolic subgroups of \(G\) of the same types as
\(P,Q\), respectively.  The pair \((P',Q')\) is in transversal position
if and only if, étale-locally, it is conjugate to \((P,Q)\).
\textup{(In §5 we shall see that the étale topology can be replaced by
the Zariski topology.)}
\item The canonical morphism \(P\times_SQ\to G\) induces a smooth
surjective morphism
\[
 P\times_SQ\longrightarrow\Gen(Q/P)
\]
and an isomorphism
\[
 (P\times_SQ)/(P\cap Q)\xrightarrow{\sim}\Gen(Q/P),
\]
where \(R=P\cap Q\) acts on \(P\times_SQ\) by
\[
 (p,q)r=(pr,r^{-1}q).
\]
\item The canonical morphism
\[
 P\longrightarrow\Par_{t(Q)}(G),\qquad p\longmapsto\Int(p)Q,
\]
induces a smooth surjective morphism
\[
 P\longrightarrow\Gen(/P)\cap\Par_{t(Q)}(G)
\]
and an isomorphism
\[
 P/(P\cap Q)\xrightarrow{\sim}
 \Gen(/P)\cap\Par_{t(Q)}(G).
\]
\item The canonical morphism
\[
 G\longrightarrow\Par_{t(P)}(G)\times_S\Par_{t(Q)}(G),\qquad
 g\longmapsto\bigl(\Int(g)P,\Int(g)Q\bigr),
\]
induces a smooth surjective morphism
\[
 G\longrightarrow
 \Gen(G)\cap
 \bigl(\Par_{t(P)}(G)\times_S\Par_{t(Q)}(G)\bigr)
\]
and an isomorphism
\[
 G/(P\cap Q)\xrightarrow{\sim}
 \Gen(G)\cap
 \bigl(\Par_{t(P)}(G)\times_S\Par_{t(Q)}(G)\bigr).
\]
\end{enumerate}
\end{sgastatement}
\SGARefTarget{sga3:xxvi:cor:4:2:5}

\SGAProof
Let us prove \textup{(i)}.  The stated condition is plainly sufficient.
Conversely, suppose \((P',Q')\) is in transversal position.  Since
\(P\) and \(P'\) are conjugate étale-locally, we may assume \(P=P'\).
It is then enough to show that two parabolic subgroups \(Q,Q'\), of the
same type and both in transversal position relative to \(P\), are
étale-locally conjugate by a section of \(P\).

By \SGARefLink{sga3:xxvi:theorem:4-2-1}{4.2.1}\textup{\SGARefLink{sga3:xxvi:theorem:4-2-1:item:vi}{(vi)}}, after another étale localization there are Borel
subgroups \(B,B'\subset P\), \(B_1\subset Q\), and \(B'_1\subset Q'\)
such that
\[
 B\cap B_1=T,\qquad B'\cap B'_1=T'
\]
are maximal tori of \(G\).  The Killing pairs \((B,T)\) and
\((B',T')\) of \(P\) are étale-locally conjugate in \(P\) (\SGARefLink{sga3:xxvi:lemma:1:16}{1.16}).
We may therefore assume \(B=B'\) and \(T=T'\).  Then
\(B_1=B'_1\) by Exp.\ XXII, \SGARefLink{sga3:xxii:prop:5.9.2}{5.9.2}, and hence \(Q=Q'\) by \SGARefLink{sga3:xxvi:lemma:3:8}{3.8}.

Parts \textup{(ii)}, \textup{(iii)}, and \textup{(iv)} are proved in
parallel.  We prove \textup{(ii)}.  Let \(g\in\Gen(Q/P)(S)\); thus
\(g\in G(S)\) and \(\Int(g)Q\) is in transversal position relative to
\(P\).  The preceding argument shows that \(Q\) and \(\Int(g)Q\) are,
étale-locally, conjugate by a section of \(P\).  Working locally for
this topology, choose \(p\in P(S)\) such that
\(\Int(g)Q=\Int(p)Q\).  Then
\[
 p^{-1}g\in\Norm_G(Q)(S)=Q(S),
\]
so \(g=pq\) for some \(q\in Q(S)\).  Thus the morphism in
\textup{\SGARefLink{sga3:xxvi:theorem:4-2-1:item:ii}{(ii)}} is an étale covering.  Comparing with
\SGARefLink{sga3:xxvi:theorem:4-2-1}{4.2.1}\textup{\SGARefLink{sga3:xxvi:theorem:4-2-1:item:ii}{(ii)}}, it is smooth and surjective.

Finally, the equivalence relation on \(P\times_SQ\) defined by
\(P\times_SQ\to G\) is exactly the relation associated with the action
of \(R=P\cap Q\),
\[
 (p,q)r=(pr,r^{-1}q).
\]
This proves the quotient assertion, since a smooth surjective morphism
is, in particular, an effective epimorphism.
\SGAQED

\begin{sgastatement}{Remark 4.2.6}
If \(P,Q\) are in transversal position, the open subset \(\Gen(Q/P)\)
of \(G\) will often be denoted \(P\cdot Q\); this notation is justified
by \SGARefLink{sga3:xxvi:corollary:4-2-5}{4.2.5}\textup{(ii)}.
\end{sgastatement}
\SGARefTarget{sga3:xxvi:remark:4:2:6}

\begin{sgastatement}{\SGAInlineRefTarget{sga3:xxvi:proposition:4-2-7}Proposition 4.2.7}
Let \(S\) be a scheme, \(G\) an \(S\)-reductive group, and \(P,Q\) two
parabolic subgroups of \(G\), with \((P,Q)\) in transversal position.
\begin{enumerate}[label=\textup{(\roman*)}]
\item The group \(P\cap Q\) is smooth over \(S\) (and in fact has
connected fibers by \SGARefLink{sga3:xxvi:remark:4:1:2}{4.1.2}).  Its identity component
\((P\cap Q)^0\) \textup{(Exp.\ VI\textsubscript{B}, \SGARefLink{thm:sga3-VIB-3.10}{3.10})} is a
subgroup of type \textup{(RC)} of \(G\)
\textup{(Exp.\ XXII, \SGARefLink{sga3:xxii:def:5.11.1}{5.11.1})}, and its unipotent radical
\textup{(loc.\ cit., 5.11.4)} decomposes as the direct product
\[
 \rad^u\bigl((P\cap Q)^0\bigr)
 =
 \bigl(\rad^u(P)\cap Q\bigr)
 \times_S
 \bigl(P\cap\rad^u(Q)\bigr).
\]
\item If \(S\) is affine, then
\[
 H^1\bigl(S,\rad^u((P\cap Q)^0)\bigr)=0.
\]
If \(S\) is semilocal, \(P\cap Q\) contains a maximal torus of \(G\).
\end{enumerate}
\end{sgastatement}
\SGARefTarget{sga3:xxvi:prop:4:2:7}

\SGAProof
The group \(P\cap Q\) is smooth by \SGARefLink{sga3:xxvi:theorem:4-2-1}{4.2.1}\textup{\SGARefLink{sga3:xxvi:theorem:4-2-1:item:ii}{(ii)}} and the
cartesian diagram
\[
\begin{tikzcd}[column sep=large,row sep=large]
 P\times_SQ \arrow[r] & G\\
 P\cap Q \arrow[u] \arrow[r] & S\arrow[u]
\end{tikzcd}.
\]

The remaining assertions about \((P\cap Q)^0\) may be checked
étale-locally.  By \SGARefLink{sga3:xxvi:theorem:4-2-1}{4.2.1}\textup{\SGARefLink{sga3:xxvi:theorem:4-2-1:item:vii}{(vii)}}, choose a splitting
\((G,T,M,R)\) of \(G\) such that \(P,Q\) contain \(T\) and are defined
by subsets \(R_1,R_2\) of \(R\), where \(R_1\) contains a system
\(R^+\) of positive roots and \(R_2\) contains the opposite system
\(R^-\).  Let \(\Delta\) be the set of simple roots of \(R^+\), and put
\[
 A_1=\Delta\cap(-R_1),\qquad
 A_2=\Delta\cap R_2,\qquad
 A=A_1\cap A_2.
\]
By \SGARefLink{sga3:xxvi:prop:1:3}{1.4}\textup{(v)} and \SGARefLink{sga3:xxvi:proposition:1-12}{1.12},
\[
\begin{aligned}
 R_1&=R^+\cup\bigl(R^-\cap(-\mathbb N A_1)\bigr),&
 R_2&=\bigl(R^+\cap\mathbb N A_2\bigr)\cup R^-,\\
 \rad^u(P)&=
   \prod_{\substack{\alpha\in R_1\\\alpha\notin-R_1}}U_\alpha,&
 \rad^u(Q)&=
   \prod_{\substack{\alpha\in R_2\\\alpha\notin-R_2}}U_\alpha.
\end{aligned}
\]
Exp.\ XXII, \SGARefLink{sga3:xxii:cor:5.6.7}{5.6.7} therefore gives
\[
\begin{aligned}
 \rad^u(P)\cap Q
  &=
  \prod_{\substack{\alpha\in R_1\cap R_2\\\alpha\notin-R_1}}U_\alpha
  =\prod_{\alpha\in K_2}U_\alpha,\\
 \rad^u(Q)\cap P
  &=
  \prod_{\substack{\alpha\in R_1\cap R_2\\\alpha\notin-R_2}}U_\alpha
  =\prod_{\alpha\in K_1}U_\alpha,
\end{aligned}
\]
where \(K_2\) consists of the positive roots that are linear
combinations of elements of \(A_2\), but not of elements of \(A\), and
\(K_1\) consists of the negative roots that are linear combinations of
elements of \(A_1\), but not of elements of \(A\).  If
\(\alpha\in K_2\) and \(\beta\in K_1\), then \(\alpha+\beta\) is neither
a root nor zero.  Hence the two groups displayed above commute.

By Exp.\ XXII, \SGARefLink{sga3:xxii:prop:5.4.5}{5.4.5}, the group
\[
 H=(P\cap Q)^0
\]
is defined by the set of roots
\[
 R_1\cap R_2
 =
 \bigl(R^+\cap\mathbb N A_2\bigr)
 \cup
 \bigl(R^-\cap(-\mathbb N A_1)\bigr).
\]
This set is closed, so \(H\) is of type \textup{(RC)} by definition
(Exp.\ XXII, \SGARefLink{sga3:xxii:def:5.11.1}{5.11.1}).  By loc.\ cit., \SGARefLink{sga3:xxii:lem:5.11.3}{5.11.3} and \SGARefLink{sga3:xxii:prop:5.11.4}{5.11.4},
\[
 \rad^u(H)=\prod_{\alpha\in K}U_\alpha,
\]
where \(K\) consists of the \(\alpha\in R_1\cap R_2\) such that
\(\alpha\notin-(R_1\cap R_2)\).  The symmetric part of
\(R_1\cap R_2\) is \(R\cap\mathbb ZA\), and therefore
\(K=K_1\cup K_2\).  This proves \textup{(i)}.

The first assertion of \textup{(ii)} follows from \textup{(i)} and
\SGARefLink{sga3:xxvi:cor:2:10}{2.10}.  For the second, the reductive group
\[
 (P\cap Q)^0/\rad^u((P\cap Q)^0)
\]
has a maximal torus \(\bar T\) when the base is semilocal.  The inverse
image \(N\) of \(\bar T\) in \((P\cap Q)^0\) is a subgroup of type
\textup{(R)} of \(G\), with solvable fibers, and
\[
 N^u=\rad^u((P\cap Q)^0)
\]
(Exp.\ XXII, \SGARefLink{sga3:xxii:prop:5.6.9}{5.6.9}).  The scheme of maximal tori of \(N\) is a
principal homogeneous bundle under \(N^u\) (loc.\ cit., \SGARefLink{sga3:xxii:cor:5.6.13}{5.6.13}), and
therefore has a section because \(H^1(S,N^u)=0\).
\SGAQED
